\SeventhArticle{3}
Nun geht unser Integral $\int y\,\mathrm{d}x$ in $\Delta\int y\,\mathrm{d}t$ über, von $t=0$ bis $t=1$ genommen. An dessen Stelle verwenden wir nach dem soeben Bemerkten $\Delta\int Y\,\mathrm{d}t$. Entwickelt man also $T^{(\mu)}$ in
\[
\alpha t^n+\beta t^{n-1}+\gamma t^{n-2}+\delta t^{n-3}+\text{usw.}
\]
so ergibt sich für $\int T^{(\mu)}\,\mathrm{d}t$, von $t=0$ bis $t=1$,
\SourceLeaf{176}
\[
=\frac{\alpha}{n+1}+\frac{\beta}{n}+\frac{\gamma}{n-1}+\frac{\delta}{n-2}+\text{usw.}
\]
Setzt man diese Größe $=M^{(\mu)}R^{(\mu)}$, so ergibt sich für das gesuchte Integral
\[
=\Delta(AR+A'R'+A''R''+A'''R'''+\text{usw.}+A^{(n)}R^{(n)})
\]
Als Beispiel geben wir die Berechnung des Koeffizienten $R''$ für $n=5$. Hier ergibt sich
\[
\begin{aligned}
T''&=5^5t^5-13\, .\,5^4t^4+59\, .\,5^3t^3-107\, .\,5^2tt+60\, .\,5\, .\,t\\[5pt]
M''&=2\times1\times(-1)\times(-2)\times(-3)=-12
\end{aligned}
\]
Daher ist $-12R''=\allowbreak \frac{3125}{6}-1625+\allowbreak \frac{7375}{4}-\frac{2675}{3}+\allowbreak 150=\allowbreak -\frac{25}{12}$, und somit $R''=\frac{25}{144}$.

Die Rechnung wird etwas kürzer, wenn $2t-1=u$ gesetzt wird. Dann ergibt sich
\[
T^{(\mu)}=\frac{\begin{gathered}(nu+n)(nu+n-2)(nu+n-4)\ldotp\ldotp\ldotp\\(nu-n+4)(nu-n+2)(nu-n)\end{gathered}}{2^n(nu+n-2\mu)}
\]
Wir setzen
\[
\frac{\begin{gathered}(nnuu-nn)\, .\,(nnuu-(n-2)^2)\, .\,(nnuu-(n-4)^2)\\\, .\,(nnuu-(n-6)^2)\ldotp\ldotp\ldotp\ldotp\end{gathered}}{nnuu-(n-2\mu)^2}=U^{(\mu)}
\]
wobei der Zähler in $\ldotp\ldotp\ldotp(nnuu-9)(nnuu-1)$ enden muss, wenn $n$ ungerade ist, oder in $\ldotp\ldotp\ldotp(nnuu-4)nu$, wenn $n$ gerade ist; dann ist
\[
T^{(\mu)}=\frac{(nu-n+2\mu)U^{(\mu)}}{2^n}
\]
Nun ist das Integral $\int T^{(\mu)}\,\mathrm{d}t$, von $t=0$ bis $t=1$ genommen, gleich dem Integral
\[
\int\tfrac12 T^{(\mu)}\,\mathrm{d}u=\int\frac{nuU^{(\mu)}\,\mathrm{d}u}{2^{n+1}}+\int\frac{(2\mu-n)U^{(\mu)}\,\mathrm{d}u}{2^{n+1}}
\]
von $u=-1$ bis $u=+1$.

Setzt man also
\[
U^{(\mu)}=\alpha u^{n-1}+\beta u^{n-3}+\gamma u^{n-5}+\delta u^{n-7}+\text{usw.}
\]
\begingroup\rightskip=0pt plus1fil\relax\noindent (denn es ist von selbst klar, dass die Potenzen $u^{n-2},\allowbreak u^{n-4},\allowbreak u^{n-6}$ usw. fehlen), so verschwindet der Teil $\int\dfrac{nuU^{(n)}\,\mathrm{d}u}{2^{n+1}}$ des Integrals bei einem ungeraden Wert von $n$, der andere Teil $\int\dfrac{(2\mu-n)U^{(\mu)}\,\mathrm{d}u}{2^{n+1}}$ dagegen bei einem geraden Wert. Daher erhält man für das Integral $\int T^{(\mu)}\,\mathrm{d}t$ bei geradem $n$\par\endgroup
\[
=\frac{n}{2^n}\left(\frac{\alpha}{n+1}+\frac{\beta}{n-1}+\frac{\gamma}{n-3}+\frac{\delta}{n-5}+\text{usw.}\right)
\]
\SourceLeaf{177}
bei ungeradem $n$ dagegen
\[
=\frac{2\mu-n}{2^n}\left(\frac{\alpha}{n}+\frac{\beta}{n-2}+\frac{\gamma}{n-4}+\frac{\delta}{n-6}+\text{usw.}\right)
\]
\Needspace{8\baselineskip}In unserem Beispiel ergibt sich
\[
U''=(25uu-25)(25uu-9)=625u^4-850uu+225,\quad\text{daher}
\]
\[
-12R''=-\tfrac1{32}\left(125-\tfrac{850}{3}+225\right)=-\tfrac{25}{12}\quad\text{wie oben.}
\]
Es ist zu beachten: $U^{(n-\mu)}=U^{(\mu)}$ und daher $\int T^{(n-\mu)}\,\mathrm{d}t=\pm\int T^{(\mu)}\,\mathrm{d}t$. Das obere Zeichen gilt bei geradem $n$, das untere bei ungeradem. Da ferner leicht einzusehen ist, dass ebenso $M^{(n-\mu)}=\pm M^{(\mu)}$ gilt, gilt stets $R^{(n-\mu)}=R^{(\mu)}$; das heißt, unter den Koeffizienten $R,R',R''\ldotp\ldotp\ldotp\ldotp R^{(n)}$ ist der letzte gleich dem ersten, der vorletzte gleich dem zweiten und so weiter.

\SeventhArticle{4}
Die von \textsc{Cotes} bis $n=10$ berechneten numerischen Werte dieser Koeffizienten schreiben wir hier aus der \textit{Harmonia Mensurarum} ab.

Für $n=1$, also bei zwei Gliedern.
\[
R=R'=\tfrac12
\]
Für $n=2$, also bei drei Gliedern.
\[
R=R''=\tfrac16,\ R'=\tfrac23
\]
Für $n=3$, also bei vier Gliedern.
\[
R=R'''=\tfrac18,\ R'=R''=\tfrac38
\]
Für $n=4$, also bei fünf Gliedern.
\[
\begin{gathered}
R=R''''=\tfrac7{90},\quad R'=R'''=\tfrac{16}{45},\\[5pt]
R''=\tfrac2{15}
\end{gathered}
\]
Für $n=5$, also bei sechs Gliedern.
\[
\begin{gathered}
R=R^{\mathrm{V}}=\tfrac{19}{288},\quad R'=R''''=\tfrac{25}{96},\\[5pt]
R''=R'''=\tfrac{25}{144}
\end{gathered}
\]
Für $n=6$, also bei sieben Gliedern.
\[
\begin{gathered}
R=R^{\mathrm{VI}}=\tfrac{41}{840},\quad R'=R^{\mathrm{V}}=\tfrac9{35},\\[5pt]
R''=R^{\mathrm{IV}}=\tfrac9{280},\quad R'''=\tfrac{34}{105}
\end{gathered}
\]
Für $n=7$, also bei acht Gliedern.
\[
\begin{gathered}
R=R^{\mathrm{VII}}=\tfrac{751}{17280},\quad R'=R^{\mathrm{VI}}=\tfrac{3577}{17280},\\[5pt]
R''=R^{\mathrm{V}}=\tfrac{49}{640},\quad R'''=R''''=\tfrac{2989}{17280}
\end{gathered}
\]
Für $n=8$, also bei neun Gliedern.
\[
\begin{gathered}
R=R^{\mathrm{VIII}}=\tfrac{989}{28350},\quad R'=R^{\mathrm{VII}}=\tfrac{2944}{14175},\\[5pt]
R''=R^{\mathrm{VI}}=-\tfrac{464}{14175},\quad R'''=R^{\mathrm{V}}=\tfrac{5248}{14175},\\[5pt]
R^{\mathrm{IV}}=-\tfrac{454}{2835}
\end{gathered}
\]
\SourceLeaf{178}
Für $n=9$, also bei zehn Gliedern.
\[
\begin{gathered}
R=R^{\mathrm{IX}}=\frac{2857}{89600},\quad R'=R^{\mathrm{VIII}}=\frac{15741}{89600},\\[5pt]
R''=R^{\mathrm{VII}}=\frac{27}{2240},\quad R'''=R^{\mathrm{VI}}=\frac{1209}{5600},\\[5pt]
R^{\mathrm{IV}}=R^{\mathrm{V}}=\frac{2889}{44800}
\end{gathered}
\]
Für $n=10$, also bei elf Gliedern.
\[
\begin{gathered}
R=R^{\mathrm{X}}=\frac{16067}{598752},\quad R'=R^{\mathrm{IX}}=\frac{26575}{149688}.\quad
R''=R^{\mathrm{VIII}}=-\frac{16175}{199584},\\[5pt]
R'''=R^{\mathrm{VII}}=\frac{5675}{12474},\quad R^{\mathrm{IV}}=R^{\mathrm{VI}}=-\frac{4825}{11088},\\[5pt]
R^{\mathrm{V}}=\frac{17807}{24948}.
\end{gathered}
\]

